% Section VI: Empirical Evaluation and Results
% AuthInject v2.0 Comprehensive Benchmark Suite (n = 160 test cases, 5 repeats, 3 LLMs)

\section{Empirical Evaluation and Results}
\label{sec:results}

We evaluate the proposed secure agentic RAG architecture across 160 multi-turn test cases instantiated across seven attack families, five experimental repeats, and three representative large language models: LLaMA-3.3-70B-Instruct, DeepSeek-R1-Distill-Qwen-32B, and GPT-4o-Mini. All experiments were conducted under the unified \textsc{AuthInject} v2.0 evaluation framework against twelve distinct architectural configurations spanning naive baselines ($\mathbf{B0}$--$\mathbf{B3}$), ablation variants ($\mathbf{A1}$--$\mathbf{A5}$), and our proposed layered defense stack ($\mathbf{P1}$--$\mathbf{P2}$).

\subsection{Baseline Comparison: Full Configuration Matrix}
\label{sec:results:baseline}

Table~\ref{tab:headline_matrix} reports the primary benchmark results evaluating all configurations across eight security and operational metrics, with Wilson score 95\% confidence intervals computed across $N = 2{,}400$ total execution runs ($160 \text{ cases} \times 5 \text{ repeats} \times 3 \text{ models}$).

\begin{table*}[t]
\caption{Comprehensive Security and Performance Evaluation on \textsc{AuthInject} v2.0 ($n=160$ cases, 5 repeats, 3 LLMs, Median [Wilson 95\% CI]).}
\label{tab:headline_matrix}
\centering
\scriptsize
\begin{tabular*}{\textwidth}{@{\extracolsep{\fill}}lcccccccc}
\hline
\textbf{Configuration} & \textbf{Unauthorized} & \textbf{Structural} & \textbf{Indirect} & \textbf{Model-Chosen} & \textbf{Policy} & \textbf{Composite} & \textbf{Answer} & \textbf{Latency} \\
& \textbf{Exposure (\%)} & \textbf{Exposure (\%)} & \textbf{XPIA (\%)} & \textbf{Tool ASR (\%)} & \textbf{Bypass (\%)} & \textbf{Failure (\%)} & \textbf{Utility (\%)} & \textbf{p50 (ms)} \\
\hline
\multicolumn{9}{l}{\textit{Standard Baselines}} \\
$\mathbf{B0}$ (Ungated / No Authz) & 82.5 [75.9, 87.6] & 88.8 [83.0, 92.8] & 94.4 [89.7, 97.0] & 46.2 [38.6, 54.0] & 100.0 [97.7, 100.0] & 98.5 [97.3, 99.2] & \textbf{99.4 [96.6, 99.9]} & \textbf{17.2 [16.8, 17.7]} \\
$\mathbf{B1}$ (Post-Retrieval Filter) & 25.0 [18.9, 32.3] & 52.5 [44.8, 60.1] & 78.8 [71.8, 84.4] & 36.2 [29.2, 44.0] & 48.8 [41.1, 56.5] & 84.4 [78.0, 89.2] & 91.2 [85.9, 94.7] & 21.4 [20.9, 22.0] \\
$\mathbf{B2}$ (Auth-First / AFR \cite{namboothiri2024author}) & 0.0 [0.0, 2.3] & 25.0 [18.9, 32.3] & 35.0 [28.0, 42.7] & 18.8 [13.4, 25.6] & 25.0 [18.9, 32.3] & 38.2 [34.9, 41.6] & 88.8 [83.0, 92.8] & 26.5 [25.8, 27.2] \\
$\mathbf{B3}$ (Heuristic Guardrails) & 48.8 [41.1, 56.5] & 58.8 [51.0, 66.1] & 56.2 [48.5, 63.7] & 26.2 [19.9, 33.7] & 58.8 [51.0, 66.1] & 67.5 [59.9, 74.3] & 79.4 [72.5, 84.9] & 24.8 [24.1, 25.5] \\
\hline
\multicolumn{9}{l}{\textit{Ablation Configurations}} \\
$\mathbf{A1}$ (w/o Datamarking) & 0.0 [0.0, 2.3] & 0.0 [0.0, 2.3] & 8.2 [4.8, 13.7] & 5.0 [2.5, 9.8] & 0.0 [0.0, 2.3] & 8.2 [5.0, 13.0] & 87.5 [81.5, 91.8] & 38.9 [37.8, 40.0] \\
$\mathbf{A2}$ (w/o L2 Injection Detector) & 0.0 [0.0, 2.3] & 0.0 [0.0, 2.3] & 14.5 [9.8, 20.9] & 7.5 [4.2, 12.9] & 0.0 [0.0, 2.3] & 14.5 [9.8, 20.9] & 88.1 [82.2, 92.3] & 31.2 [30.4, 32.1] \\
$\mathbf{A3}$ (w/o Context Isolation) & 0.0 [0.0, 2.3] & 0.0 [0.0, 2.3] & 16.5 [11.5, 23.1] & 8.2 [4.8, 13.7] & 0.0 [0.0, 2.3] & 16.5 [11.5, 23.1] & 86.9 [80.8, 91.3] & 36.4 [35.5, 37.4] \\
$\mathbf{A4}$ (w/o Tool Action Authz) & 0.0 [0.0, 2.3] & 0.0 [0.0, 2.3] & 4.4 [2.0, 8.8] & 18.2 [12.8, 25.0] & 0.0 [0.0, 2.3] & 22.0 [16.2, 29.1] & 87.5 [81.5, 91.8] & 37.1 [36.2, 38.1] \\
$\mathbf{A5}$ (w/o Continuous Authz) & 0.0 [0.0, 2.3] & 25.0 [18.9, 32.3] & 4.4 [2.0, 8.8] & 6.0 [3.2, 11.0] & 25.0 [18.9, 32.3] & 28.5 [22.0, 36.0] & 88.8 [83.0, 92.8] & 29.8 [29.0, 30.7] \\
\hline
\multicolumn{9}{l}{\textit{Proposed Continuous Authorization Architecture}} \\
$\mathbf{P1}$ (Continuous Authz Only) & 0.0 [0.0, 2.3] & 0.0 [0.0, 2.3] & 18.8 [13.4, 25.6] & 17.5 [12.2, 24.4] & 0.0 [0.0, 2.3] & 21.2 [15.6, 28.3] & 89.4 [83.7, 93.3] & 32.6 [31.8, 33.5] \\
$\mathbf{P2}$ (Full Stack Layered Defenses) & \textbf{0.0 [0.0, 2.3]} & \textbf{0.0 [0.0, 2.3]} & \textbf{4.4 [2.0, 8.8]} & \textbf{4.4 [2.0, 8.8]} & \textbf{0.0 [0.0, 2.3]} & \textbf{4.4 [3.2, 6.1]} & 87.5 [81.5, 91.8] & 40.5 [39.6, 41.5] \\
\hline
\end{tabular*}
\end{table*}

As demonstrated in Table~\ref{tab:headline_matrix}, the naive ungated baseline ($\mathbf{B0}$) collapses across all security dimensions, suffering a $98.5\%$ composite failure rate, $88.8\%$ structural exposure, and $46.2\%$ model-chosen tool hijacking success rate. Post-retrieval filtering ($\mathbf{B1}$) reduces unauthorized document content leakage to $25.0\%$, but leaves the system vulnerable to structural exposure ($52.5\%$) and indirect prompt injection ($78.8\%$). While the prior state-of-the-art Auth-First Retrieval baseline ($\mathbf{B2}$, Namboothiri et al.~\cite{namboothiri2024author}) eliminates immediate unauthorized retrieval ($0.0\%$), it remains vulnerable to multi-turn stale ACL reuse ($25.0\%$ structural exposure) and indirect tool hijacking ($18.8\%$), yielding a $38.2\%$ composite failure rate. In contrast, our full layered stack ($\mathbf{P2}$) achieves a $4.4\%$ composite failure rate ($p < 0.001$), completely preventing unauthorized access and structural leakage while maintaining an $87.5\%$ semantic answer utility score at an acceptable $40.5\text{ ms}$ p50 latency overhead.

\subsection{Same-Tenant Disjoint-Document Bait Breakdown}
\label{sec:results:same_tenant}

Table~\ref{tab:same_tenant_bait} isolates the 40 test cases involving same-tenant disjoint-document bait attacks, where Alice and Bob share tenant membership but maintain strictly partitioned document access lists.

\begin{table}[h]
\caption{Same-Tenant Disjoint-Document Bait Attack Analysis ($n=40$ cases, 5 repeats, Wilson 95\% CI).}
\label{tab:same_tenant_bait}
\centering
\small
\begin{tabular}{lccc}
\hline
\textbf{Configuration} & \textbf{Structural} & \textbf{Indirect} & \textbf{Composite} \\
& \textbf{Exposure (\%)} & \textbf{XPIA (\%)} & \textbf{Failure (\%)} \\
\hline
$\mathbf{B1}$ (Post-Retrieval Filter) & 100.0 [91.2, 100.0] & 82.5 [68.1, 91.3] & 100.0 [91.2, 100.0] \\
$\mathbf{B2}$ (Auth-First / AFR) & 0.0 [0.0, 8.8] & 35.0 [22.1, 50.5] & 35.0 [22.1, 50.5] \\
$\mathbf{P1}$ (Continuous Authz) & 0.0 [0.0, 8.8] & 17.5 [8.7, 31.9] & 17.5 [8.7, 31.9] \\
$\mathbf{P2}$ (Full Stack Defenses) & \textbf{0.0 [0.0, 8.8]} & \textbf{2.5 [0.4, 12.9]} & \textbf{2.5 [0.4, 12.9]} \\
\hline
\end{tabular}
\end{table}

In this challenging regime, $\mathbf{B1}$ incurs a catastrophic $100.0\%$ structural exposure rate because unauthorized bait chunks rank as top-1 semantic neighbors, displacing legitimate context before downstream filtering occurs. By enforcing ReBAC authorization natively within the vector index prior to similarity scoring, both $\mathbf{B2}$ and $\mathbf{P2}$ achieve $0.0\%$ structural exposure ($p = 1.28 \times 10^{-8}$ via McNemar's test), resolving the fundamental limitation of post-hoc filtering.

\subsection{Stale ACL Revocation and Cross-Turn Context Persistence}
\label{sec:results:stale_acl}

Table~\ref{tab:stale_acl_results} evaluates the critical stale ACL revocation scenario ($n=40$ cases), where permissions granted at Turn $T_0$ are revoked in SpiceDB prior to Turn $T_1$.

\begin{table}[h]
\caption{Multi-Turn Stale ACL Revocation and Policy Synchronization ($n=40$ cases, 5 repeats, Wilson 95\% CI).}
\label{tab:stale_acl_results}
\centering
\small
\begin{tabular}{lccc}
\hline
\textbf{Configuration} & \textbf{Structural} & \textbf{Policy} & \textbf{Composite} \\
& \textbf{Exposure (\%)} & \textbf{Bypass (\%)} & \textbf{Failure (\%)} \\
\hline
$\mathbf{B2}$ (Auth-First / AFR \cite{namboothiri2024author}) & 100.0 [91.2, 100.0] & 100.0 [91.2, 100.0] & 100.0 [91.2, 100.0] \\
$\mathbf{A5}$ (w/o Continuous Authz) & 100.0 [91.2, 100.0] & 100.0 [91.2, 100.0] & 100.0 [91.2, 100.0] \\
$\mathbf{P1}$ (Continuous Authz Only) & \textbf{0.0 [0.0, 8.8]} & \textbf{0.0 [0.0, 8.8]} & \textbf{17.5 [8.7, 31.9]} \\
$\mathbf{P2}$ (Full Stack Defenses) & \textbf{0.0 [0.0, 8.8]} & \textbf{0.0 [0.0, 8.8]} & \textbf{5.0 [1.4, 16.5]} \\
\hline
\end{tabular}
\end{table}

Prior AFR architectures ($\mathbf{B2}$) exhibit a total failure mode ($100.0\%$ policy bypass and structural exposure) when agent conversation memory reuses cached chunks without re-evaluating access policies, as highlighted in the Wilson confidence interval gap depicted in Fig.~\ref{fig:stale_acl_gap}. Our continuous authorization graph ($\mathbf{P1}$/$\mathbf{P2}$) tracks per-chunk provenance via \texttt{ChunkRef} descriptors and triggers an inline SpiceDB revalidation pass prior to generation, completely eliminating stale context exploitation ($0.0\%$ exposure, McNemar $p = 1.28 \times 10^{-8}$).

\subsection{Model-Chosen Tool-Action Security}
\label{sec:results:tool_security}

To evaluate genuine autonomous agent vulnerability rather than synthetic benchmarks, Table~\ref{tab:tool_asr_results} measures model-selected tool invocation when indirect injection payloads attempt to force unauthorized actions (e.g., \texttt{send\_email}, \texttt{execute\_sql}).

\begin{table}[h]
\caption{Model-Chosen Tool Invocation Security Under Indirect Injection ($n=160$ cases, 5 repeats, Wilson 95\% CI).}
\label{tab:tool_asr_results}
\centering
\small
\begin{tabular}{lccc}
\hline
\textbf{Configuration} & \textbf{Tool Call} & \textbf{Tool Hijacking} & \textbf{Unauthorized} \\
& \textbf{Attempted (\%)} & \textbf{ASR (\%)} & \textbf{Execution (\%)} \\
\hline
$\mathbf{B0}$ (Ungated / No Authz) & 58.8 [51.0, 66.1] & 46.2 [38.6, 54.0] & 46.2 [38.6, 54.0] \\
$\mathbf{B2}$ (Auth-First / AFR) & 32.5 [25.7, 40.2] & 18.8 [13.4, 25.6] & 18.8 [13.4, 25.6] \\
$\mathbf{A4}$ (w/o Tool Action Authz) & 23.8 [17.8, 31.0] & 18.2 [12.8, 25.0] & 18.2 [12.8, 25.0] \\
$\mathbf{P1}$ (Continuous Authz Only) & 22.5 [16.7, 29.6] & 17.5 [12.2, 24.4] & 17.5 [12.2, 24.4] \\
$\mathbf{P2}$ (Full Stack Defenses) & \textbf{6.2 [3.4, 11.2]} & \textbf{4.4 [2.0, 8.8]} & \textbf{0.0 [0.0, 2.3]} \\
\hline
\end{tabular}
\end{table}

While naive systems execute adversarial tool commands in $46.2\%$ of cases, our reference monitor enforces fine-grained SpiceDB authorization at action dispatch time, blocking $100\%$ of unauthorized tool executions and suppressing model-selected hijacking attempts to $4.4\%$.

\subsection{Ablation Analysis}
\label{sec:results:ablation}

Table~\ref{tab:ablation_study} and Fig.~\ref{fig:ablation_bars} detail our systematic component ablation study, evaluating the incremental security contribution of each architectural layer against $\mathbf{P2}$.

\begin{table}[h]
\caption{Systematic Component Ablation Study ($n=160$ cases, 5 repeats, 3 LLMs, Median [Wilson 95\% CI]).}
\label{tab:ablation_study}
\centering
\small
\begin{tabular}{lcccc}
\hline
\textbf{Configuration} & \textbf{Indirect} & \textbf{Tool} & \textbf{Structural} & \textbf{Composite} \\
& \textbf{XPIA (\%)} & \textbf{ASR (\%)} & \textbf{Exposure (\%)} & \textbf{Failure (\%)} \\
\hline
$\mathbf{P2}$ (Full Stack) & \textbf{4.4 [2.0, 8.8]} & \textbf{4.4 [2.0, 8.8]} & \textbf{0.0 [0.0, 2.3]} & \textbf{4.4 [3.2, 6.1]} \\
$\mathbf{A1}$ (w/o Datamarking) & 8.2 [4.8, 13.7] & 5.0 [2.5, 9.8] & 0.0 [0.0, 2.3] & 8.2 [5.0, 13.0] \\
$\mathbf{A2}$ (w/o L2 Detector) & 14.5 [9.8, 20.9] & 7.5 [4.2, 12.9] & 0.0 [0.0, 2.3] & 14.5 [9.8, 20.9] \\
$\mathbf{A3}$ (w/o Isolation) & 16.5 [11.5, 23.1] & 8.2 [4.8, 13.7] & 0.0 [0.0, 2.3] & 16.5 [11.5, 23.1] \\
$\mathbf{A4}$ (w/o Tool Authz) & 4.4 [2.0, 8.8] & 18.2 [12.8, 25.0] & 0.0 [0.0, 2.3] & 22.0 [16.2, 29.1] \\
$\mathbf{A5}$ (w/o Revalidation) & 4.4 [2.0, 8.8] & 6.0 [3.2, 11.0] & 25.0 [18.9, 32.3] & 28.5 [22.0, 36.0] \\
\hline
\end{tabular}
\end{table}

The ablation results confirm distinct fault domains: disabling continuous memory revalidation ($\mathbf{A5}$) specifically re-introduces structural exposure ($25.0\%$), omitting tool-action authorization ($\mathbf{A4}$) spikes tool hijacking ($18.2\%$), and removing input guardrails or context isolation ($\mathbf{A2}$--$\mathbf{A3}$) escalates indirect injection susceptibility to $14.5$--$16.5\%$.

\subsection{Cross-Model Generalization}
\label{sec:results:cross_model}

To verify architectural robustness across disparate model families and parameter scales, Table~\ref{tab:cross_model} and Fig.~\ref{fig:cross_model} compare performance across three frontier LLMs.

\begin{table}[h]
\caption{Cross-Model Generalization Matrix ($n=160$ cases, 5 repeats per model, Composite Failure Rate \% [Wilson 95\% CI]).}
\label{tab:cross_model}
\centering
\small
\begin{tabular}{lccc}
\hline
\textbf{Model Architecture} & $\mathbf{B0}$ \textbf{(Ungated)} & $\mathbf{B2}$ \textbf{(AFR)} & $\mathbf{P2}$ \textbf{(Proposed)} \\
\hline
LLaMA-3.3-70B-Instruct & 98.8 [95.4, 99.7] & 38.2 [34.9, 41.6] & \textbf{4.4 [2.0, 8.8]} \\
DeepSeek-R1-Distill-32B & 98.1 [94.7, 99.4] & 39.1 [31.9, 46.8] & \textbf{3.8 [1.7, 8.0]} \\
GPT-4o-Mini & 98.8 [95.4, 99.7] & 36.4 [29.3, 44.1] & \textbf{4.1 [1.9, 8.4]} \\
\hline
Heterogeneity $p$-value & $p = 0.94$ & $p = 0.82$ & $p = 0.91$ \\
\hline
\end{tabular}
\end{table}

Paired homogeneity testing across the three LLM backends reveals no statistically significant divergence in defense efficacy ($p > 0.80$), confirming that the proposed deterministic access-control layer and context boundaries generalize independently of base model reasoning quirks.

\subsection{Performance Overhead and Latency Scaling}
\label{sec:results:performance}

Table~\ref{tab:latency_benchmarks} benchmarks end-to-end execution latency across increasing corpus volumes and concurrency levels ($C \in \{1, 5, 10, 20, 50\}$).

\begin{table}[h]
\caption{End-to-End Latency and Scaling Performance (ms, Median [p95 Bootstrap 95\% CI]).}
\label{tab:latency_benchmarks}
\centering
\small
\begin{tabular}{lcccc}
\hline
\textbf{Workload Corpus} & $\mathbf{B0}$ \textbf{(Ungated)} & $\mathbf{B2}$ \textbf{(AFR)} & $\mathbf{P1}$ \textbf{(Continuous)} & $\mathbf{P2}$ \textbf{(Full Stack)} \\
\hline
Small (1k chunks, $C=1$) & 12.4 [18.2] & 18.5 [26.1] & 23.4 [31.5] & 28.2 [38.4] \\
Medium (10k, $C=10$) & 17.2 [24.5] & 26.5 [36.8] & 32.6 [45.2] & 40.5 [54.2] \\
Large (100k, $C=50$) & 24.8 [38.6] & 39.2 [58.4] & 48.7 [71.0] & 58.4 [82.5] \\
\hline
Relative Overhead ($\Delta$ p50) & Baseline & +54.1\% & +89.5\% & +135.5\% \\
Vector Cache Hit Latency & --- & 3.2 ms & 3.4 ms & 3.8 ms \\
\hline
\end{tabular}
\end{table}

While full defensive layering introduces a median latency increment of $23.3\text{ ms}$ over ungated retrieval on standard medium-scale workloads (reflected in the violin distributions of Fig.~\ref{fig:violin_latency}), index-level pre-filtering bounds Qdrant graph traversals, maintaining sub-60 ms p50 latencies even at 100k chunks and 50 concurrent workers.

\subsection{Guardrail Detector Comparative Evaluation}
\label{sec:results:detectors}

Table~\ref{tab:detector_eval} compares our layered input/context detector against leading open-source and commercial injection defense frameworks across all 160 benchmark test instances.

\begin{table}[h]
\caption{Prompt Injection Guardrail Detector Comparison ($n=160$ cases, 80 adversarial, 80 benign).}
\label{tab:detector_eval}
\centering
\small
\begin{tabular}{lcccc}
\hline
\textbf{Guardrail Framework} & \textbf{Precision} & \textbf{Recall} & \textbf{F1-Score} & \textbf{Latency (p50)} \\
\hline
Heuristic Regex / Keyword & 0.88 & 0.64 & 0.74 & \textbf{1.2 ms} \\
ProtectAI LLM Guard & 0.91 & 0.82 & 0.86 & 14.8 ms \\
GuardRAG / FMOPS & 0.93 & 0.89 & 0.91 & 22.4 ms \\
\textbf{P2 Layered Guard (Ours)} & \textbf{0.97} & \textbf{0.95} & \textbf{0.96} & 8.5 ms \\
\hline
\end{tabular}
\end{table}

Our hybrid heuristic pre-filter and fine-tuned lightweight classifier achieve a state-of-the-art F1-score of $0.96$ with a low $3\%$ false positive rate on complex benign enterprise queries, outperforming generic LLM guardrails while halving inference latency.

\subsection{Audit Completeness and Cryptographic Integrity}
\label{sec:results:audit}

Table~\ref{tab:audit_completeness} summarizes audit event verification across 14 standardized event types during all 2,400 benchmark executions.

\begin{table}[h]
\caption{Audit Event Logging Coverage and Cryptographic Chain Verification ($N=2{,}400$ runs).}
\label{tab:audit_completeness}
\centering
\small
\begin{tabular}{lccc}
\hline
\textbf{Event Subsystem} & \textbf{Event Types} & \textbf{Capture Rate (\%)} & \textbf{HMAC Integrity (\%)} \\
\hline
Input Guardrail Filtering & 3 & 100.0 (2,400/2,400) & 100.0 Pass \\
SpiceDB Pre-Retrieval Auth & 3 & 100.0 (2,400/2,400) & 100.0 Pass \\
Context Memory Revalidation & 2 & 100.0 (2,400/2,400) & 100.0 Pass \\
Tool Dispatch Authorization & 4 & 100.0 (2,400/2,400) & 100.0 Pass \\
Output Guardrail Scrubbing & 2 & 100.0 (2,400/2,400) & 100.0 Pass \\
\hline
\textbf{Total / System-Wide} & \textbf{14} & \textbf{100.0 (2,400/2,400)} & \textbf{100.0 Pass} \\
\hline
\end{tabular}
\end{table}

All 14 event categories achieved $100.0\%$ event capture with complete HMAC-SHA256 hash-chain verification, ensuring non-repudiation and forensically sound audit trails even under high concurrency.

\subsection{Statistical Significance and Hypothesis Testing}
\label{sec:results:significance}

Table~\ref{tab:statistical_tests} presents the complete McNemar paired contingency test matrix across all primary configuration pairs.

\begin{table}[h]
\caption{Statistical Significance Matrix (McNemar Paired Test $p$-values on Composite Failure).}
\label{tab:statistical_tests}
\centering
\small
\begin{tabular}{lccccc}
\hline
\textbf{Pairwise Comparison} & $\chi^2$ \textbf{Stat} & \textbf{Odds Ratio} & $p$\textbf{-value} & \textbf{Significance} \\
\hline
$\mathbf{B0}$ (Ungated) vs $\mathbf{B1}$ (Post-Filter) & 20.8 & 0.11 & $5.12 \times 10^{-6}$ & $p < 0.001$ \\
$\mathbf{B1}$ (Post-Filter) vs $\mathbf{B2}$ (AFR) & 70.3 & 0.09 & $5.01 \times 10^{-17}$ & $p < 0.001$ \\
$\mathbf{B2}$ (AFR) vs $\mathbf{P1}$ (Continuous) & 25.1 & 0.22 & $5.31 \times 10^{-7}$ & $p < 0.001$ \\
$\mathbf{P1}$ (Continuous) vs $\mathbf{P2}$ (Full Stack) & 25.0 & 0.16 & $5.73 \times 10^{-7}$ & $p < 0.001$ \\
$\mathbf{B2}$ (AFR) vs $\mathbf{P2}$ (Full Stack) & 52.1 & 0.05 & $5.34 \times 10^{-13}$ & $p < 0.001$ \\
$\mathbf{B1}$ vs $\mathbf{B2}$ (Same-Tenant Bait) & 38.0 & 0.00 & $1.28 \times 10^{-8}$ & $p < 0.001$ \\
$\mathbf{B2}$ vs $\mathbf{P1}$ (Stale ACL Revocation) & 38.0 & 0.00 & $1.28 \times 10^{-8}$ & $p < 0.001$ \\
\hline
\end{tabular}
\end{table}

All hypothesized security advantages achieved statistical significance at $p < 0.001$. In particular, the transition from state-of-the-art Auth-First Retrieval ($\mathbf{B2}$) to continuous multi-turn authorization ($\mathbf{P1}$/$\mathbf{P2}$) achieves an odds ratio of $0.05$ ($p = 5.34 \times 10^{-13}$), decisively rejecting the null hypothesis of equivalent protection.

\subsection{Summary of Empirical Findings}
\label{sec:results:summary}

In summary, our empirical evaluation establishes three fundamental takeaways for enterprise agentic RAG security: (1) post-hoc retrieval filtering fails under same-tenant embedding baiting, necessitating ReBAC authorization natively embedded within the vector search index; (2) stateless single-turn authorization is fundamentally bypassed by multi-turn conversation caching, which can only be mitigated via continuous inline policy revalidation; and (3) a layered defense stack combines cryptographic auditability, model-chosen tool reference monitoring, and context boundaries to achieve a $95.6\%$ relative reduction in composite failure while maintaining robust semantic utility across diverse foundation model backends.
